\scanpage{71}
% archive foliation: 70
% lines 1--2
Paradigme des courbes algébriques

I\quad Point de vue topologique.
% lines 3--7
On considère des surfaces topologiques (sans bord)
orientables paracompactes \inserted{$X$} -- on se donne
 aussi un ensemble $\omega$ de cardinal $2$,
et on pose
\[T=Z\wedge_{Z/2}\omega,\]
% lines 8--14
et on suppose donnée sur $X$ une $\omega$-orientation,
i.e. une \inserted{$Z/2$-}application
\deleted{\unsure{équivariante|compatible}}\inserted{de $\omega$ dans}
l'ens. des orientations de $X$, ou ce
qui revient au $=$, un isomorphisme
du faisceau des entiers tordus sur $X$
avec le faisceau constant $T_X$.
% lines 15--22
On suppose que tte composante connexe
$X_i$ de $X$ \deleted{soit}\inserted{est} \unsure{homéomorphes|homéomorphe} à un
$\widehat X_i\setminus S_i$, où $\widehat X_i$ est une surface
compacte, et $S_i$ une partie finie de
$\widehat X_i$. On sait que $(\widehat X_i,S_i)$ est déterminé
à isom. unique près par $X_i$
(comme le \unsure{décomp.} des compactifiés de Čech
